\section{Homomorphic / Holographic Consensus DAO}
\label{sec:dao}

\subsection{Two adjectives, two distinct properties}

The \Zoo{} \DAO{} is described as both \emph{homomorphic} and
\emph{holographic}. They are independent properties:

\begin{itemize}[nosep,leftmargin=1.2em]
\item \textbf{Homomorphic} means votes are tallied while still encrypted.
  Tallying happens under fully homomorphic encryption (\FHE{})
  on \Lux{} F-Chain (\Lux{} \textsc{lp-013}~\cite{lp-013-fchain}).
  The aggregate result is decryptable; individual ballots are not.
\item \textbf{Holographic} means every part of the \DAO{} encodes the
  whole's intent. Nested sub-\DAO{}s (per-\textsc{llm} chains, per
  research vertical, per regional chapter) inherit governance
  topology from the meta-\DAO{} but can also propose and ratify
  locally. Each shard is structurally a small copy of the whole.
\end{itemize}

\subsection{Homomorphic vote tallying}

A vote is encoded as a ciphertext under an F-Chain
\FHE{} public key. The protocol:

\begin{enumerate}[nosep,leftmargin=1.4em]
\item Each participant submits $\mathrm{Enc}_{\text{pk}}(v_i \cdot w_i)$
  where $v_i \in \{-1, 0, +1\}$ is the vote (against / abstain / for)
  and $w_i$ is the participant's voting weight (\S\ref{sec:weighted}).
\item F-Chain~\cite{lp-013-fchain} computes $\sum_i \mathrm{Enc}(v_i w_i)$
  homomorphically. The encrypted sum is published on-chain.
\item Decryption of the aggregate is performed only by the
  M-Chain MPC custodial threshold (\Lux{} \textsc{lp-019}~\cite{lp-019-mchain})
  after the voting window closes.
\item Individual ballots remain encrypted forever. Aggregate result
  is published with a Quasar 3.0 finality proof.
\end{enumerate}

This eliminates voter coercion: a participant cannot prove how they
voted to a coercer because nobody --- not even the M-Chain custodians
acting jointly --- can decrypt an individual ciphertext. Only the sum
is decryptable.

\subsection{Holographic structure}

The meta-\DAO{} is the apex. Sub-\DAO{}s exist for:

\begin{itemize}[nosep,leftmargin=1.2em]
\item Each per-\textsc{llm} chain (one sub-\DAO{} per Zen model;
  see~\cite{zoo-per-llm-chains}).
\item Each research vertical (Conservation AI, Educational AI,
  Frontier AI, DeSci).
\item Each regional chapter (geographic conservation cells).
\end{itemize}

Each sub-\DAO{} inherits the meta-\DAO{}'s governance topology
(weighted voting, \FHE{} tally, M-Chain custody) and can propose and
ratify locally on issues scoped to its shard. Issues that cross shards
escalate to the meta-\DAO{}. The structural property is
\emph{self-similarity}: a sub-\DAO{} looks like a smaller meta-\DAO{}
in every governance dimension.

This is what we mean by \emph{holographic}: every part contains the
whole's pattern. A sub-\DAO{} can fork into its own chain with the same
governance machinery and remain interoperable with the meta-\DAO{}
without bespoke integration.

\subsection{Weighted voting}
\label{sec:weighted}

Voting weight is computed as a weighted sum across four axes:

\begin{itemize}[nosep,leftmargin=1.2em]
\item \textbf{Advocacy} --- public posts, citations, and social
  engagement, attested through Z-Chain (anti-Sybil; see below).
\item \textbf{Involvement} --- contributor activity: commits, code
  reviews, training submissions, ceremony participation, validator
  uptime.
\item \textbf{Contribution} --- landmark research artifacts, model
  releases, novel architectures, dataset curation, reproducibility
  attestations.
\item \textbf{Token stake} --- locked $\$$ZOO stake. Always present
  but deliberately the smallest weight, to prevent plutocracy.
\end{itemize}

The formula:

\begin{equation}
\text{weight}(p) \;=\; \alpha \cdot a_p \;+\; \beta \cdot i_p \;+\; \gamma \cdot c_p \;+\; \delta \cdot s_p
\end{equation}

\noindent subject to $\alpha + \beta + \gamma + \delta = 1$, where
$a_p, i_p, c_p, s_p$ are the participant's normalized advocacy,
involvement, contribution, and token-stake scores respectively.

Default initial weights at activation (2025-12-25) are governed by the
meta-\DAO{} and start as:

\[
(\alpha, \beta, \gamma, \delta) \;=\; (0.20,\ 0.30,\ 0.40,\ 0.10)
\]

\noindent i.e.\ 20\% advocacy, 30\% involvement, 40\% contribution,
10\% token stake. These weights can be adjusted by meta-\DAO{} vote
(itself a homomorphic, weighted vote).

\subsection{Anti-Sybil}

Each input score is gated by attestation:

\begin{itemize}[nosep,leftmargin=1.2em]
\item Advocacy ($a_p$) requires the social platform to be
  attestation-rooted (i.e.\ the platform itself runs an attestation
  endpoint anchored on Z-Chain). Posts on un-rooted platforms do not
  count.
\item Involvement ($i_p$) is computed from the contribution-graph
  attestation pipeline anchored on \Lux{} A-Chain
  (\Lux{} \textsc{lp-134}~\cite{lp-134-topology}). Each commit and
  review is signed by a \textsc{tee}-attested identity.
\item Contribution ($c_p$) requires reproducibility attestation
  (per ZIP-0606 in the Zoo \textsc{zip} corpus); a paper, model, or
  dataset must verify before it confers contribution weight.
\item Token stake ($s_p$) is unforgeable on-chain by definition but
  is rate-limited by the $\delta = 0.10$ default to bound plutocracy.
\end{itemize}

A purely token-rich attacker can buy at most $\delta = 0.10$ of total
voting weight. To dominate, the attacker must also generate provable
advocacy, involvement, and contribution --- which the attestation
pipeline rejects unless backed by real activity.

\subsection{Threshold custody of governance keys}

Governance keys (the master key for upgrade proposals, treasury
movements above policy thresholds, and emergency halts) are held in
threshold custody on M-Chain (\Lux{} \textsc{lp-019}~\cite{lp-019-mchain}).
A successful proposal pulls a threshold partial signature from each
custodian; the resulting Corona \textsc{pq} aggregate signature
authorizes the on-chain action. No single custodian can act
unilaterally; no custodian can be compelled by physical coercion of
one party.

\subsection{Cross-references}

\begin{itemize}[nosep,leftmargin=1.2em]
\item \Lux{} \textsc{lp-013} --- F-Chain FHE substrate (vote
  tally)~\cite{lp-013-fchain}.
\item \Lux{} \textsc{lp-019} --- M-Chain MPC (governance key
  custody)~\cite{lp-019-mchain}.
\item \Lux{} \textsc{lp-134} --- A/B/M/F topology (contributor
  attestation on A-Chain)~\cite{lp-134-topology}.
\item \textsc{zip-0017} --- DAO Governance Framework (Zoo \textsc{zip}
  corpus).
\item \textsc{zip-0026} --- Ecosystem Reputation System (the
  involvement-score input).
\end{itemize}
